% Shared bibliography (\input by literature_review.tex and manuscript.tex). Every entry verified against its
% published abstract or proceedings page.
\begin{thebibliography}{99}\small
\bibitem{foulds2020bayes} J.~R. Foulds, R.~Islam, K.~N. Keya, S.~Pan. Bayesian modeling of intersectional fairness: the variance of bias. \emph{Proc.\ SIAM Int.\ Conf.\ Data Mining (SDM)}, 424--432, 2020. doi:10.1137/1.9781611976236.48. arXiv:1811.07255.
\bibitem{foulds2020icde} J.~R. Foulds, R.~Islam, K.~N. Keya, S.~Pan. An intersectional definition of fairness. \emph{Proc.\ IEEE 36th Int.\ Conf.\ Data Engineering (ICDE)}, 1918--1921, 2020. doi:10.1109/ICDE48307.2020.00203.
\bibitem{kearns2018} M.~Kearns, S.~Neel, A.~Roth, Z.~S. Wu. Preventing fairness gerrymandering: auditing and learning for subgroup fairness. \emph{Proc.\ ICML}, PMLR 80:2564--2572, 2018.
\bibitem{hebert2018} \'U.~H\'ebert-Johnson, M.~P. Kim, O.~Reingold, G.~N. Rothblum. Multicalibration: calibration for the (computationally-identifiable) masses. \emph{Proc.\ ICML}, PMLR 80:1939--1948, 2018.
\bibitem{molina2022} M.~Molina, P.~Loiseau. Bounding and approximating intersectional fairness through marginal fairness. \emph{Advances in NeurIPS 35}, 2022. arXiv:2206.05828.
\bibitem{wang2022} A.~Wang, V.~V. Ramaswamy, O.~Russakovsky. Towards intersectionality in machine learning: including more identities, handling underrepresentation, and performing evaluation. \emph{Proc.\ ACM FAccT}, 336--349, 2022. doi:10.1145/3531146.3533101.
\bibitem{gohar2023} U.~Gohar, L.~Cheng. A survey on intersectional fairness in machine learning: notions, mitigation, and challenges. \emph{Proc.\ IJCAI} (Survey Track), 6619--6627, 2023. doi:10.24963/ijcai.2023/742.
\bibitem{ji2020} D.~Ji, P.~Smyth, M.~Steyvers. Can I trust my fairness metric? Assessing fairness with unlabeled data and Bayesian inference. \emph{Advances in NeurIPS 33}, 18600--18612, 2020. arXiv:2010.09851.
\bibitem{morina2019} G.~Morina, V.~Oliinyk, J.~Waton, I.~Maru\v{s}i\'c, K.~Georgatzis. Auditing and achieving intersectional fairness in classification problems. arXiv:1911.01468, 2019.
\bibitem{lum2022} K.~Lum, Y.~Zhang, A.~Bower. De-biasing ``bias'' measurement. \emph{Proc.\ ACM FAccT}, 379--389, 2022. doi:10.1145/3531146.3533105.
\bibitem{herlihy2024} C.~Herlihy, K.~Truong, A.~Chouldechova, M.~Dud\'ik. A structured regression approach for evaluating model performance across intersectional subgroups. \emph{Proc.\ ACM FAccT}, 2024. arXiv:2401.14893.
\bibitem{himmelreich2025} J.~Himmelreich, A.~Hsu, E.~Veomett, K.~Lum. The intersectionality problem for algorithmic fairness. \emph{Algorithmic Fairness Through the Lens of Metrics and Evaluation}, PMLR 279:68--95, 2025. arXiv:2411.02569.
\bibitem{cherian2024} J.~J. Cherian, E.~J. Cand\`es. Statistical inference for fairness auditing. \emph{Journal of Machine Learning Research} 25(149):1--49, 2024.
\bibitem{hardt2016} M.~Hardt, E.~Price, N.~Srebro. Equality of opportunity in supervised learning. \emph{Advances in NIPS 29}, 3315--3323, 2016.
\bibitem{cotter2019} A.~Cotter, M.~Gupta, H.~Jiang, N.~Srebro, K.~Sridharan, S.~Wang, B.~Woodworth, S.~You. Training well-generalizing classifiers for fairness metrics and other data-dependent constraints. \emph{Proc.\ ICML}, PMLR 97:1397--1405, 2019.
\bibitem{ding2021} F.~Ding, M.~Hardt, J.~Miller, L.~Schmidt. Retiring Adult: new datasets for fair machine learning. \emph{Advances in NeurIPS 34}, 6478--6490, 2021. arXiv:2108.04884.
\bibitem{folktables} Folktables project. Dataset definitions, documentation and data-access tools. \url{https://github.com/socialfoundations/folktables}.
\bibitem{fogliato2020} R.~Fogliato, A.~Chouldechova, M.~G'Sell. Fairness evaluation in presence of biased noisy labels. \emph{Proc.\ AISTATS}, PMLR 108:2325--2336, 2020.
\bibitem{berger1986} J.~O. Berger, L.~M. Berliner. Robust Bayes and empirical Bayes analysis with $\epsilon$-contaminated priors. \emph{Annals of Statistics} 14(2):461--486, 1986. doi:10.1214/aos/1176349933.
\bibitem{datta1995} G.~S. Datta, P.~Lahiri. Robust hierarchical Bayes estimation of small area characteristics in the presence of covariates and outliers. \emph{Journal of Multivariate Analysis} 54(2):310--328, 1995.
\end{thebibliography}
